\section{Fields of definition and canonical bases}
\label{sec:fod}

Canonical bases of types are, in algebraically closed fields, the smallest fields of
definition of loci.  We construct them by linear algebra (reduced row echelon
bases), which is the most economical route for a formalization: it needs no
Gr\"obner bases and no Galois descent.

\subsection{Subspaces of a coordinate space}

Fix a finite set $\mathcal M$ (the index set of coordinates; later $\mathcal M$ is a set
of monomials) and consider $K$-subspaces $W\subseteq K^{\mathcal M}$.  Fix a total order
on $\mathcal M$.

\begin{definition}[defined over a subfield]\label{def:subspace-defined}
Let $F\subseteq K$ be a subfield.  $W$ is \emph{defined over $F$} if $W$ is spanned
by $W\cap F^{\mathcal M}$; equivalently (by comparing dimensions over $F$ and $K$) if
$W=(W\cap F^{\mathcal M})\otimes_F K$ under the natural map.
\end{definition}

\begin{definition}[reduced row echelon basis]\label{def:rref}
A finite sequence $w_1,\dots,w_r$ of vectors of $K^{\mathcal M}$ is in \emph{reduced
row echelon form} if, writing $\mu_i\in\mathcal M$ for the first index (in the fixed
order) at which $w_i$ is nonzero (the \emph{pivot}), we have $\mu_1<\mu_2<\dots<\mu_r$,
$w_i(\mu_i)=1$ and $w_j(\mu_i)=0$ for all $j\ne i$.
\end{definition}

\begin{lemma}[existence and uniqueness]\label{lem:rref}
Every subspace $W\subseteq K^{\mathcal M}$ has exactly one basis in reduced row echelon
form, the \emph{RREF basis} $\operatorname{rref}(W)$.  For an automorphism $\sigma$ of the
field $K$, $\operatorname{rref}(\sigma W)=\sigma(\operatorname{rref}(W))$.  The same holds
with $K$ replaced by any field.
\end{lemma}
\begin{proof}
Existence: Gaussian elimination on any basis.  Uniqueness: let $(w_i)$ and $(w'_j)$ be
two RREF bases.  The set of pivots of an RREF basis of $W$ is determined by $W$: it is
the set of $\mu\in\mathcal M$ such that some vector of $W$ has pivot $\mu$ (if $v\in W$
has pivot $\mu$, write $v=\sum c_i w_i$; the first $i$ with $c_i\neq0$ has
$\mu_i=\mu$; conversely each $w_i$ has pivot $\mu_i$).  So both bases have the same
pivots $\mu_1<\dots<\mu_r$.  Then $w_i-w'_i\in W$ vanishes at every pivot, so it is
zero (a nonzero vector of $W$ has a pivot).  Equivariance: $\sigma$ preserves the
property ``in reduced row echelon form''.
\end{proof}

\begin{definition}[field of definition of a subspace]\label{def:fod-subspace}
$\FOD(W)$ is the subfield of $K$ generated by the prime field and all coordinates
of all vectors of $\operatorname{rref}(W)$.
\end{definition}

\begin{lemma}\label{lem:fod-subspace}
Let $F\subseteq K$ be a subfield.  $W$ is defined over $F$ iff $\FOD(W)\subseteq F$.
Hence $\FOD(W)$ is the smallest field of definition of $W$, it is finitely generated
over the prime field, and $\sigma(\FOD(W))=\FOD(\sigma W)$ for every automorphism
$\sigma$ of $K$.  In particular, for $\sigma\in\Aut(K)$: $\sigma W=W$ iff $\sigma$ fixes
$\FOD(W)$ pointwise.
\end{lemma}
\begin{proof}
If $\FOD(W)\subseteq F$, the RREF basis lies in $F^{\mathcal M}$, so $W$ is defined over
$F$.  Conversely, if $W$ is spanned by $W\cap F^{\mathcal M}$, Gaussian elimination over
$F$ of an $F$-basis of $W\cap F^{\mathcal M}$ produces an RREF sequence in
$F^{\mathcal M}$ which spans $W$ over $K$; by uniqueness it is $\operatorname{rref}(W)$,
so $\FOD(W)\subseteq F$.  Equivariance is \cref{lem:rref}.  If $\sigma W=W$ then $\sigma$
maps $\operatorname{rref}(W)$ to $\operatorname{rref}(W)$ as a sequence, hence fixes each
coordinate; the converse is clear.
\end{proof}

\subsection{Ideals}

Let $X=(X_1,\dots,X_n)$ and, for $d\ge0$, let $K[X]_{\le d}$ be the subspace of
polynomials of total degree $\le d$, identified with $K^{\mathcal M_d}$ where
$\mathcal M_d$ is the set of monomials of degree $\le d$ (ordered, say,
lexicographically).  For an ideal $I\subseteq K[X]$ put $I_{\le d}=I\cap K[X]_{\le d}$.

\begin{lemma}\label{lem:ideal-filtration}
Let $F\subseteq K$ be a subfield and $I\subseteq K[X]$ an ideal.
\begin{enumerate}[nosep]
\item $I$ is defined over $F$ (\cref{def:locus}) iff every $I_{\le d}$ is defined over
  $F$ as a subspace of $K[X]_{\le d}$.
\item If $I_{\le d}$ is defined over $F$ and $d'\le d$ then $I_{\le d'}$ is defined over
  $F$; hence $\FOD(I_{\le d'})\subseteq\FOD(I_{\le d})$.
\item There is $d_0$ such that $I$ is generated by $I_{\le d_0}$; for $d\ge d_0$,
  $I$ is defined over $F$ iff $I_{\le d}$ is defined over $F$.
\end{enumerate}
\end{lemma}
\begin{proof}
(1) ``$\Rightarrow$'': $I=J\cdot K[X]$ with $J=I\cap F[X]$.  Choosing an $F$-basis of
$K$ containing $1$ (\cref{lem:extcontr}), $I=\bigoplus_\beta J\beta$ and the degree
filtration is compatible with this decomposition, so $I_{\le d}=\bigoplus_\beta J_{\le d}\beta$
is spanned by $J_{\le d}\subseteq I_{\le d}\cap F[X]$.  ``$\Leftarrow$'': $I=\bigcup_d I_{\le d}$
is generated by $\bigcup_d(I_{\le d}\cap F[X])\subseteq I\cap F[X]$.
(2) $I_{\le d'}=I_{\le d}\cap K[X]_{\le d'}$, and the intersection of two subspaces
defined over $F$ is defined over $F$: tensoring the exact sequence
$0\to U\cap U'\to U\oplus U'\to V$ of $F$-spaces with the flat $F$-module $K$ shows
$(U\otimes K)\cap(U'\otimes K)=(U\cap U')\otimes K$.  Then apply \cref{lem:fod-subspace}.
Here $\FOD(I_{\le d'})$ is computed in $K^{\mathcal M_{d'}}$ and $\FOD(I_{\le d})$ in
$K^{\mathcal M_d}$; since $\FOD$ of a subspace is its smallest field of definition
(\cref{lem:fod-subspace}), it does not depend on the ambient coordinate space in which
the subspace is placed (a subspace of $K^{\mathcal M_{d'}}$ is defined over $F$ iff its
image in $K^{\mathcal M_d}$ under the coordinate inclusion is), nor on the chosen
monomial order.
(3) $K[X]$ is noetherian; if $I$ is generated by polynomials of degree $\le d_0$ and
$I_{\le d}$ ($d\ge d_0$) is defined over $F$ then $I$ is generated by elements of
$I_{\le d}\cap F[X]$, so $I$ is defined over $F$; the converse is (1).
\end{proof}

\begin{definition}[smallest field of definition]\label{def:fod}
For an ideal $I\subseteq K[X]$, $\FOD(I):=\bigcup_{d}\FOD(I_{\le d})=\FOD(I_{\le d_0})$,
an increasing union which stabilises by \cref{lem:ideal-filtration}.
\end{definition}

\begin{theorem}[Weil]\label{thm:weil-fod}
Let $I\subseteq K[X]$ be an ideal.  A subfield $F\subseteq K$ is a field of definition
of $I$ iff $\FOD(I)\subseteq F$.  $\FOD(I)$ is finitely generated over the prime field.
For every automorphism $\sigma$ of $K$: $\FOD(\sigma I)=\sigma\FOD(I)$, and $\sigma I=I$
iff $\sigma$ fixes $\FOD(I)$ pointwise.
\end{theorem}
\begin{proof}
Combine \cref{lem:ideal-filtration} with \cref{lem:fod-subspace}, noting
$(\sigma I)_{\le d}=\sigma(I_{\le d})$.
\end{proof}

\begin{remark}
The same construction applies verbatim to ideals of $L[X]$ for any field $L$ with
subfields $F\subseteq L$, and for $J\subseteq L[X]$ one has $\FOD(J\cdot K[X])=\FOD(J)$
because $\operatorname{rref}((J\cdot K[X])_{\le d})=\operatorname{rref}(J_{\le d})$
(the RREF basis of $J_{\le d}$ over $L$ is in reduced row echelon form and spans
$(J K[X])_{\le d}$ by \cref{lem:ideal-filtration}(1)).  We shall only use ideals of
$K[X]$.
\end{remark}

\begin{leannote}
Mathlib has the predicate \lean{Matrix.IsReducedRowEchelon} but neither existence nor
uniqueness of the reduced row echelon basis; \cref{lem:rref} must be formalized (it is
elementary).  \lean{Submodule.restrictScalars}, \lean{Submodule.span} and flatness of
field extensions (\lean{Module.Flat}) supply the facts used in
\cref{lem:ideal-filtration}.  An alternative is to define $\FOD$ via Mathlib's
\lean{MonomialOrder} division algorithm and reduced Gr\"obner bases, but uniqueness of
reduced Gr\"obner bases is not in Mathlib either; the RREF route is shorter.
\end{leannote}

\subsection{Canonical bases}

\begin{definition}[canonical base]\label{def:cb}
Let $a$ be a tuple and $L$ a small algebraically closed subfield.  The
\emph{canonical base} of $a$ over $L$ is the subfield
\[
   \Cb(a/L):=k\cdot\FOD(P_L(a))\subseteq L,
\]
the field generated by $k$ and the smallest field of definition of the locus ideal
$P_L(a)$ (which lies in $L$ because $P_L(a)$ is defined over $L$).  It is a small
subfield.  For a finite tuple $w$ we write $\Cb(a/w)=\Cb(a/\acl(w))$.
\end{definition}

\begin{proposition}[properties of canonical bases]\label{prop:cb}
Let $a,a'$ be tuples, $L,L'$ small algebraically closed subfields and $C$ a small
algebraically closed subfield with $C\subseteq L$.
\begin{enumerate}[label=\textup{(Cb\arabic*)}]
\item\label{Cb:ind} $\rk(a/\Cb(a/L))=\rk(a/L)$, i.e.\ $a\ind{\Cb(a/L)}L$.
\item\label{Cb:char} $a\ind{C}L$ iff $\Cb(a/L)\subseteq C$.
\item\label{Cb:equiv} $\sigma(\Cb(a/L))=\Cb(\sigma a/\sigma L)$ for $\sigma\in\Gamma$; and for
  $\sigma\in\Gamma_L$, $\sigma$ fixes $\Cb(a/L)$ pointwise iff $P_L(\sigma a)=P_L(a)$ iff
  $\sigma a$ and $a$ are $\Gamma_L$-conjugate.
\item\label{Cb:two} If $a\ind{L}L'$ and $a\ind{L'}L$ then $\Cb(a/L)=\Cb(a/L')=\Cb(a/\acl(LL'))$.
\item\label{Cb:interalg} If $\acl(a)=\acl(a')$ then $\acl(\Cb(a/L))=\acl(\Cb(a'/L))$.
\item\label{Cb:basechange} If $e$ is a tuple with $e\ind{}aL$ then
  $\Cb(a/\acl(Le))=\Cb(a/L)$.
\end{enumerate}
\end{proposition}
\begin{proof}
\ref{Cb:ind}: Let $F=\Cb(a/L)$ and $F'=\acl(F)\subseteq L$, a small algebraically
closed field over which $P:=P_L(a)$ is defined.  Let $Q=P\cap F'[X]$, a prime of
$F'[X]$ with $Q\cdot K[X]=P$.  Since $a$ is a zero of $I_L(a)$ it is a zero of $P$ and
hence of $Q$, so $Q\subseteq I_{F'}(a)$.  By the ``consequently'' clause of \cref{lem:dim} ($\dim P$ may be computed with a
generic point over any small algebraically closed field over which $P$ is defined; a
generic point of $P$ over $F'$ has rank $\dim P=\dim P_L(a)=\rk(a/L)$, and its
vanishing ideal over $F'$ is $Q$ by \cref{lem:extcontr}) we get
$\trdeg_{F'}\operatorname{Frac}(F'[X]/Q)=\rk(a/L)$, while
$\trdeg_{F'}\operatorname{Frac}(F'[X]/I_{F'}(a))=\rk(a/F')$.  By \cref{lem:dimdrop}
(non-strict form, as $Q\subseteq I_{F'}(a)$), $\rk(a/F')\le\rk(a/L)$; the
reverse inequality is \ref{R:mono}.  Finally $\rk(a/F)=\rk(a/F')$ by \ref{R:acl}.

\ref{Cb:char}: If $a\ind{C}L$ then $P_C(a)=P_L(a)$ by \cref{lem:dim}(2), and $P_C(a)$
is defined over $C$, so $\Cb(a/L)=k\FOD(P_C(a))\subseteq C$.  Conversely, if
$\Cb(a/L)\subseteq C\subseteq L$ then by \ref{R:mono} and \ref{Cb:ind}
$\rk(a/L)\le\rk(a/C)\le\rk(a/\Cb(a/L))=\rk(a/L)$.

\ref{Cb:equiv}: $\sigma P_L(a)=P_{\sigma L}(\sigma a)$ and \cref{thm:weil-fod}.  For
$\sigma\in\Gamma_L$, $P_L(\sigma a)=\sigma P_L(a)$ equals $P_L(a)$ iff $\sigma$ fixes
$\FOD(P_L(a))$ pointwise; and $P_L(\sigma a)=P_L(a)$ iff $\sigma a$, $a$ are
$\Gamma_L$-conjugate by \cref{lem:dim}(3).

\ref{Cb:two}: Put $M=\acl(LL')$.  From $a\ind{L}L'$ we get $\rk(a/M)=\rk(a/LL')=\rk(a/L)$
by \ref{R:acl}, i.e.\ $a\ind{L}M$, so $P_L(a)=P_M(a)$ by \cref{lem:dim}(2) and
$\Cb(a/L)=\Cb(a/M)$.  Symmetrically $\Cb(a/L')=\Cb(a/M)$.

\ref{Cb:interalg}: Let $C=\acl(\Cb(a'/L))\subseteq L$.  By \ref{Cb:ind} and \ref{R:acl},
$a'\ind{C}L$; since $a\subseteq\acl(a')$, \ref{R:aclinv} gives $a\ind{C}L$, so
$\Cb(a/L)\subseteq C$ by \ref{Cb:char}, i.e.\ $\acl(\Cb(a/L))\subseteq\acl(\Cb(a'/L))$.
By symmetry the two closures coincide.

\ref{Cb:basechange}: $e\ind{}aL$ gives $e\ind{L}a$ and hence $a\ind{L}e$, i.e.
$a\ind{L}\acl(Le)$ by \ref{R:acl}; then $P_{\acl(Le)}(a)=P_L(a)$ by \cref{lem:dim}(2).
\end{proof}

\begin{remark}[what canonical bases are not]\label{rem:cb-not-acl}
In general $\Cb(a/L)\not\subseteq\acl(a)$: for independent $s,s',t$ and
$a=(t,st+s')$, $L=\acl(s,s')$, the locus of $a$ is the line $X_2=sX_1+s'$ and
$\Cb(a/L)=k(s,s')$, while $\rk(a)=2$ and $s\notin\acl(a)$.  The main theorem's
hypothesis ``$a\interalg\Cb(xy/a)$'' is therefore a genuine condition (see
\cref{sec:gc}).
\end{remark}

\begin{leannote}
$P_L(a)$ is \lean{(AclGeom.idealOf L a).map (MvPolynomial.map (algebraMap L K))}.  The
rank-two special case of \ref{Cb:char}--\ref{Cb:ind} (the coefficient closure of a plane
curve equation is the whole parameter plane) is \lean{AclGeom.QWitness.Psi.exists\_incidence\_germs\_with\_full\_coefficients}
in \lean{AclGeom/Config/ParameterField.lean}; there the field of definition of a principal
prime ideal is read off from the lexicographically monic generator
(\lean{AclGeom.FiniteCorrespondencePair.curveEquation}).  For principal ideals the two
descriptions agree: if $P=(f)$ with $f$ monic for a monomial order then
$\FOD(P)$ is the field generated by the coefficients of $f$ (the RREF basis of
$P_{\le d}$ consists of the monomial multiples of $f$ reduced against each other, whose
coefficients lie in the field generated by those of $f$, and conversely $f$ is the
unique element of $P$ with its leading monomial and coefficient $1$ reduced modulo
the others, so its coefficients lie in every field of definition).
\end{leannote}
